% ID: SEL-D0022-OB-E03
% ESTADO: cotejada. Segunda pasada visual completa y borrador compilado revisados.
% Fuente: Selectividad/matematicas_aplicadas/reserva_a_EXAMEN 5 -ANDALUCÍA.pdf
% SHA-256: 668689642ca997cf817c9f9f6608d2835efe898e3162b050d7b22b51f225d363
% Página: 2; opción B; ejercicio original 3.

De los sucesos $A$ y $B$ de un experimento aleatorio se conocen las siguientes probabilidades:
\[
P(A)=0.4,\qquad P(B)=0.5,\qquad P\bigl((A\cup B)^C\bigr)=0.1.
\]
\begin{enumerate}[label=\alph*)]
\item \textbf{(0.75 puntos)} Razone si $A$ y $B$ son sucesos compatibles.
\item \textbf{(0.75 puntos)} Razone si $A$ y $B$ son sucesos independientes.
\item \textbf{(0.5 puntos)} Calcule $P(A\cap B^C)$.
\item \textbf{(0.5 puntos)} Calcule $P(A/B^C)$.
\end{enumerate}
